\documentclass[11pt]{article}
\usepackage[margin=25mm]{geometry}
\usepackage{fontspec}\setmainfont{Latin Modern Roman}
\usepackage{amsmath,amssymb,mathtools}
\usepackage{unicode-math}\setmathfont{Latin Modern Math}
\usepackage{xurl,hyperref}
\hypersetup{colorlinks=true,urlcolor=blue}
\setcounter{secnumdepth}{0}
\setlength{\parindent}{0pt}\setlength{\parskip}{5pt}
\setlength{\emergencystretch}{3em}
\newcommand{\sref}[2]{\hyperlink{source:#1:#2}{[S#2]}}
\newcommand{\cataloguescope}[1]{\paragraph{Recorded target.}#1}
\begin{document}
\section{Kourovka Notebook Problem 21.117}\label{statement}

\subsection{Definitions and mathematical statement}

\paragraph{Shared notation.}$\mathbb N=\{1,2,\ldots\}$, and $\mathbb Z,\mathbb Q,\mathbb R,\mathbb C$ denote the integers, rationals, reals and complex numbers. Any other conventions are specified locally.


For a finite generating set $S$ of a group $G$, let $B_S(n)$ be the words of length at most $n$ in $S\cup S^{-1}$ and put $\lambda(G,S)=\lim_{n\to\infty}|B_S(n)|^{1/n}$. Exponential growth means $\lambda(G,S)>1$; uniform exponential growth means $\inf_S\lambda(G,S)>1$, with the infimum over all finite generating sets. An infinite group is simple if its only normal subgroups are $1$ and itself. It is just-infinite if every nontrivial normal subgroup has finite index, and hereditarily just-infinite if every finite-index subgroup is just-infinite.
\paragraph{Statement recorded in the supplied source.}
Does there exist (a) a finitely generated infinite simple group, and (b) a finitely generated hereditarily just-infinite group, of exponential growth but without uniform exponential growth? These are two separate existence questions. \sref{statement}{1}\sref{statement}{2}
\subsection{Short English statement}
Ask whether simple groups and hereditarily just-infinite groups can have exponential growth whose rate becomes arbitrarily close to one as the generators change.
\subsection{Sources}
\begin{description}
\item[\hypertarget{source:statement:1}{S1}] ULAM.AI and the UnsolvedMath Contributors, \emph{UnsolvedMath: A Curated Collection of Open Mathematics Problems}, catalogue entry. CC BY 4.0. \url{https://huggingface.co/datasets/ulamai/UnsolvedMath}.
\item[\hypertarget{source:statement:2}{S2}] E. I. Khukhro and V. D. Mazurov (editors), The Kourovka Notebook, Unsolved Problems in Group Theory, No.21, version48 (6 October2026), Problem 21.117 and its remarks. \url{https://arxiv.org/abs/1401.0300v48}.
\end{description}
\label{end:statement}
\end{document}
